\documentclass[11pt,a4paper]{article}
\usepackage{fontspec}
\setmainfont{Noto Serif}
\setsansfont{Noto Sans}
\usepackage{babel}
\babelprovide[main,import]{russian}
\usepackage{amsmath,amssymb,mathtools,booktabs,geometry,fancyhdr,microtype,xcolor,hyperref}
\geometry{left=21mm,right=21mm,top=21mm,bottom=21mm}
\definecolor{ink}{HTML}{173D57}
\hypersetup{colorlinks=true,urlcolor=ink,linkcolor=ink,pdftitle={BGS-F. Спектральный возврат в размерности h209}}
\pagestyle{fancy}\fancyhf{}\lhead{\small BGS-F}\rhead{\small Спектральный возврат при фиксированном кике}\cfoot{\thepage}
\setlength{\headheight}{15pt}\setlength{\parindent}{1em}\setlength{\parskip}{2.5pt}
\setlength{\emergencystretch}{3em}\tolerance=3000
\newcommand{\Tr}{\operatorname{Tr}}
\newcommand{\EE}{\mathbb E}
\newcommand{\R}{\mathbb R}
\newcommand{\Z}{\mathbb Z}
\newcommand{\C}{\mathbb C}
\newcommand{\si}{\mathrm{Sine}_2}
\newcommand{\eps}{\varepsilon}
\newcommand{\sinc}{\operatorname{sinc}}
\newcommand{\sect}[1]{\section*{\sffamily\color{ink}#1}}
\newcommand{\subt}[1]{\par\medskip\noindent\textbf{#1.}\quad}
\begin{document}
\noindent{\small\sffamily\color{ink} МАТЕМАТИЧЕСКАЯ ФИЗИКА\hfill 19 НОЯБРЯ 2026}
\begin{center}{\LARGE\sffamily Спектральный возврат для BGS-F}\\[4pt]
{\large Точный перенос по решётке и конечный сертификат при $N=h_{209}$}
\end{center}
\sect{Результат и его область}
При фиксированном $\tau=10^{-80}$ для исходного оператора $U_N$ получена строгая оценка
\[
 \boxed{\quad |\Tr U_N^{209}|/N>1.8\cdot10^{-5},\qquad
 N=h_{209}\simeq5.0013614231\cdot10^{79}.\quad}                 \tag{1}
\]
Здесь $(1+\sqrt2)^k=h_k+v_k\sqrt2$. В этой размерности $\tau N>1/2$. Для сравнения, у матрицы Хаара $\EE|\Tr V^{209}|^2=209$. Доказательство (1) сохраняет комплексные фазы. Вычислительная его часть состоит из свёрток над рациональными числами и четырёх возведений в степень над гауссовыми целыми с контролируемым округлением.
    
Из (1) получены два следствия для того же $N$. При $s=10^6$ значение нормированного абелева форм-фактора превышает предложенный предельный ориентир более чем на $3.2\cdot10^{-10}$. Для ограниченной непрерывной статистики недозаполнения окна с треугольным весом и радиусом $10^6$ разность математических ожиданий превышает $8.74\cdot10^{-6}$. Размер окна фиксирован в единицах средней щели. Оба результата относятся к полному возмущённому спектру.

Для любого нечётного пеллева индекса $k\ge3$ ниже доказано тождество переноса, из которого получен сертификат. Отдельно доказаны необходимые следствия предполагаемой слабой сходимости для всех нечётных $N$. В частности, она запрещает макроскопический нормированный след равномерно на каждой полосе $0<|m|\le CN$. При этой же гипотезе выполнение одного теста (P) уже обеспечивает сходимость всего парного семейства, включая (A).

\subt{Статус основной задачи} Доказательство или опровержение предела BGS-F остаётся открытым. Оценка (1) относится к одной размерности. Для бесконечного ряда $h_k$ не установлена равномерная положительная нижняя граница следа. Парный предел (P) по всем нечётным $N$ также не вычислен.

\subt{Знак траекторного вклада} Из (1) следует также $Y_{209,N}>3.23\cdot10^{-10}N^2$ в обозначениях постановки. В самом деле, циклически эквивалентные пары дают $D_{m,N}=\sum_x d(x)|a(x)|^2\le m\Tr P_N^m\le mN$, где $P_N(x,y)=|U_N(x,y)|^2$ стохастична и $d(x)\le m$ есть период слова. Поэтому $Y_{m,N}=|p_{m,N}|^2-D_{m,N}$ имеет заявленную нижнюю границу. Здесь положительная интерференция обнаружена при одном коротком возврате.

\subt{Содержание} На с.~2 определена модель. На с.~3 выведено тождество переноса. На с.~4--5 даны рациональная аппроксимация и бюджет ошибок. На с.~6 доказаны статистические следствия сертификата. На с.~7--8 исследована связь слабого закона с парными тестами и указана остающаяся оценка.
\newpage
\sect{1. Модель и арифметическое возвращение}
В $\ell^2(\Z/N\Z)$, $N\ge3$ нечётно, пусть $e_N(t)=\exp(2\pi i t/N)$, черта над целым означает его обратный класс по модулю $N$. При $d=(a,b)$ положим
\[
 (T(d)\psi)(x)=e_N(b(x+\bar2a))\psi(x+a),\quad
 M_{xy}=N^{-1/2}e_N[-\bar8(3x^2-2xy+3y^2)] .
\]
Пусть $r(q)=(1-\cos14\pi q)(1+\cos2\pi q)$, $e^{i\alpha_0}=(-3+4i)/5$ и
\[
 f(q,p)=\cos14\pi q+\tfrac1{10}r(q)+\tfrac1{100}r(q)
    [\sin2\pi q+\cos(2\pi p-\alpha_0)] .
\]
Если $f-1/10=\sum c_{ba}e^{2\pi i(bq+ap)}$, то $A=\sum c_{ba}T(a,b)$,
$U=\exp(i\xi A)M$, $\xi=N/10^{80}$. Реальность $f$ и $T(d)^*=T(-d)$ дают $A=A^*$. Унитарность $M$ следует из ортогональности строк дискретной матрицы Фурье, так как $\bar4$ обратим.
При $\omega((a,b),(c,d))=ad-bc$ выполняются
\[
 T(u)T(v)=e_N(\bar2\,\omega(u,v))T(u+v),\qquad
 M T(v) M^*=T(Sv),\quad S=\begin{pmatrix}3&4\\2&3\end{pmatrix}.       \tag{2}
\]
Первое тождество получается применением сдвигов к произвольной функции. Во втором сравним ядра $MT(a,b)$ и $T(a',b')M$. Коэффициенты $y$ и $x$ в разности с квадратичной фазой $M_{xy}$ равны соответственно $3a/4+b$ и $-a/4$ слева, $a'/4$ и $-3a'/4+b'$ справа. Отсюда $a'=3a+4b$, $b'=2a+3b$. Свободные члены с обеих сторон совпадают и равны $-3a^2/8-ab/2$.

Для нечётного $k$ имеем $h_k^2-2v_k^2=-1$. Умножение чисел вида $u+v\sqrt2$ соответствует умножению матриц $\left(\begin{smallmatrix}u&2v\\v&u\end{smallmatrix}\right)$. Поэтому при $N=h_k$
\[
 S^k-I=2N B_k,\qquad
 B_k=\begin{pmatrix}h_k&2v_k\\ v_k&h_k\end{pmatrix},\qquad \det B_k=-1. \tag{3}
\]
Рекурсия $h_{j+1}=h_j+2v_j$, $v_{j+1}=h_j+v_j$ из $(1,0)$ показывает, что $N$ нечётно. Из (2), (3) следует $M^k=\mu I$, $|\mu|=1$. Действительно, оператор, коммутирующий с $T(0,1)$, диагонален, а коммутирование с $T(1,0)$ делает его скаляром.

Фазы $U_N$ обозначаем $e^{2\pi i\theta_j}$. Исходный процесс есть
$\Xi_N=\sum_{j,\ell}\delta_{N(\theta_j-\alpha+\ell)}$, где $\alpha$ равномерна на $[0,1)$. Для любого интервала длины $L$ имеем $\EE\Xi_N(I)=L$ интегрированием по $\alpha$. Целевой процесс имеет ядро $K(x,y)=\sin(\pi(x-y))/(\pi(x-y))$ с диагональю $1$.
\newpage
\sect{2. Формула переноса на половинной решётке}
Возьмём любой конечный полином $V=\sum_{d\in D}z_dT(d)$ с целыми векторами $d$. Обозначим $H_N$ матрицу на $c\in\frac12\Z^2$ со стрелками
\[
 c\longmapsto c'=Sc+d,\qquad
 H_N(c',c)=z_d e_N\!\left(\bar4\,\omega(d,2Sc)\right).
\]
Пусть $E$ содержит все состояния циклов длины $k$ и конечно. Ограничение на строки и столбцы из $E$ обозначим $H_{N,E}$. Тогда
\[
 \boxed{\frac{\Tr((VM)^k)}{N\mu}=\Tr(H_{N,E}^{\,k}) .}                 \tag{4}
\]
\textit{Доказательство.} Итерируем на операторах $X\mapsto VMXM^*$, начиная с $I$. Каждое слово $d_0,\ldots,d_{k-1}$ порождает $v_0=0$, $v_{j+1}=Sv_j+d_j$ в $\Z^2$. На шаге $j$ его множитель равен $z_{d_j}e_N(\bar2\omega(d_j,Sv_j))$. При взятии нормированного следа остаются ровно слова с $v_k\in N\Z^2$, поскольку $N^{-1}\Tr T(v)=\mathbf1_{v=0\bmod N}$. Левая часть (4) получается из $M^{-k}=\bar\mu I$.

Цикл с заданным началом чтения того же слова имеет единственное вещественное начальное состояние $c_0=(I-S^k)^{-1}v_k$, а $c_j=v_j+S^jc_0$. По (3) условие $v_k\in N\Z^2$ эквивалентно $c_0\in\tfrac12\Z^2$. Все $c_j$ тогда принадлежат этой решётке. В $\Z/N\Z$ половины интерпретируются с помощью $\bar2$. Добавка к суммарной симплектической фазе есть
\[
 \sum_{j=0}^{k-1}\omega(d_j,S^{j+1}c_0)
 =\omega(c_k,S^kc_0)-\omega(c_0,c_0)=0\pmod N .
\]
Здесь использованы $d_j=c_{j+1}-Sc_j$ и $\det S=1$. Тем самым и амплитуда, и кратность каждого упорядоченного слова сохранены. Это доказывает (4).

\subt{Конечная область} В координатах $u=a+\sqrt2 b$, $v=a-\sqrt2 b$ отображение $S$ умножает $u,v$ на $\lambda=3+2\sqrt2$ и $\lambda^{-1}$. Если $D_\pm=\max_{d\in D}|d_a\pm\sqrt2 d_b|$, то на каждом цикле
\[
 |u|\le D_+/(\lambda-1),\qquad |v|\le D_-/(1-\lambda^{-1}).              \tag{5}
\]
Первое следует из обратной рекурсии в точке максимума $|u|$, второе из прямой рекурсии для $v$. Оценка не зависит от длины цикла. Поскольку $S\equiv I\pmod2$, класс $2c\bmod2$ сохраняется. Достаточно четырёх отдельных матриц.

Для произвольного $\xi$ можно взять $V=\sum_{j=0}^J(i\xi A)^j/j!$ в точной алгебре (2). Интегральный остаток и самосопряжённость дают $\|V-e^{i\xi A}\|\le |\xi|^{J+1}\|A\|^{J+1}/(J+1)!$. Телескопирование $k$ множителей даёт ошибку не более $k\eps(1+\eps)^{k-1}$, если одношаговая ошибка не превосходит $\eps$. Для (1) выгоднее следующий вариант полинома.
\newpage
\sect{3. Рациональная аппроксимация кика}
В этом и следующем разделах $k=209$, $N=h_k$, $Q=2^{61}$. Используем рациональный параметр $x=500136142312/10^{12}$; $|x-\xi|<10^{-12}$. Пишем $W(a,b)=(-1)^{ab}T(a,b)$. Тогда точный множитель в произведении $W(d)W(e)$ равен $\exp(\pi i\omega(d,e)/N)$ перед $W(d+e)$ для целых меток. В коэффициентах $A$ переход к $W$ вносит тот же знак. Это сохраняет сдвиг $q+1/2$ в физическом символе смешанного члена.

Пусть $C(a,b)=(-1)^{ab}c_{ba}$, а звёздочка означает обычную свёртку на $\Z^2$. Сначала вычислим
\[
 t_d=(-1)^{d_ad_b}\sum_{j=0}^{16}\frac{(ix)^j}{j!}(C^{*j})(d),
 \qquad C^{*0}=\delta_0 .                                                 \tag{6}
\]
Все $c_{ba}$ лежат в $\frac1{4000}\Z[i]$. Для воспроизведения, коэффициенты $4r$ при частотах $0,\pm1,\pm7,\pm6,\pm8$ соответственно равны $4,2,-2,-1,-1$. По формуле для $f$ это даёт 32 ненулевые метки и $\sum_d(|\Re c_d|+|\Im c_d|)=1.186$. Здесь и далее индекс $d$ в $c_d$ означает пару индексов $ba$, согласованную с меткой $d=(a,b)$.

Оставим метки с $|\Re t_d|+|\Im t_d|>10^{-11}$ и округлим обе компоненты на сетку $Q^{-1}\Z$. Обозначим результат $z_d$, остальные коэффициенты положим нулевыми. Получены 693 метки. Для $V=\sum z_dT(d)$ имеем
\[
 \|V-e^{i\xi A}\|\le\eps,\qquad \eps=8\cdot10^{-10}.                  \tag{7}
\]
Детали оценки важны ввиду малости искомого следа. Расхождение параметров стоит не более $1.186|x-\xi|$. Хвост ряда после члена 16 не превосходит $2(1.186 x)^{17}/17!$, по отношению последовательных членов. На слове длины $j\le16$ отброшенная фаза в базисе $W$ ограничена сверху по модулю угла числом $18\pi\binom j2/N$, поскольку $|d_a|\le1$, $|d_b|\le9$. Умножаем на сумму модулей коэффициентов каждого порядка, затем суммируем. Итог меньше $10^{-70}$. Вклад обрезки равен сверху сумме $|\Re t_d|+|\Im t_d|$ по удалённым меткам, вклад округления не больше $693/Q$. Рациональное сложение этих пяти границ даёт менее $5.546\cdot10^{-10}$.

\begin{center}\small
\begin{tabular}{@{}lr@{}}\toprule
Источник ошибки одного кика & Верхняя граница \\ \midrule
Замена параметра & $2.562\cdot10^{-13}$ \\
Хвост экспоненты & $7.833\cdot10^{-19}$ \\
Замена фаз Вейля в коэффициентах & $5.183\cdot10^{-76}$ \\
Отброшенные коэффициенты & $5.543\cdot10^{-10}$ \\
Округление оставшихся коэффициентов & $3.006\cdot10^{-16}$ \\ \bottomrule
\end{tabular}\end{center}

\subt{Дискретная область} Пусть $D_0=\max_D(|a|+1.415|b|)$, $u_0=D_0/4.82$, $v_0=D_0/0.828$, $R=\lceil u_0+v_0\rceil+1=118$. Берём целые $X,Y=2c_a,2c_b$ с $|X|,|Y|\le R$ и
\[
 |1000X+1414Y|\le\lceil2000u_0\rceil+R,\quad
 |1000X-1414Y|\le\lceil2000v_0\rceil+R .
\]
Они включают (5), поскольку $0<\sqrt2-1.414<0.001$. Размеры в порядке классов $00,01,10,11$ равны $2399,2396,2398,2396$.
\newpage
\sect{4. Сертификат следа и бюджет ошибок}
В (4) при $(A',B')=S(X,Y)$ положим $t=aB'-bA'$. Имеем
\[
 e_N(\bar4 t)=i^{-Nt}\exp(2\pi i t/(4N)).
\]
Вводим матрицу $H$ с рациональными комплексными элементами, оставляя в этой формуле множитель $i^{-Nt}$. Максимальная построчная и постолбцовая сумма модулей ошибки ограничена сверху величиной $\sum_D|z_d|\,2\pi\max|t|/(4N)$. Для каждого $d$ стрелка является частичной перестановкой. Из $|X|,|Y|\le118$, $|a|\le16$, $|b|\le144$, $\sum|z_d|<2$ следует $\|H-H_{N,E}\|<10^{-70}$ (критерий Шура).

Для оценки степеней нужна хорошая граница нормы. На всей половинной решётке определим $L_de_c=i^{-N\omega(d,2c)}e_{c+d}$. Тогда $L_dL_e=(-1)^{\omega(d,e)}L_{d+e}$, $L_d^*=L_{-d}$. Следовательно, $(-1)^{d_ad_b}L_d$ образуют обычное унитарное представление $\Z^2$. Формула (6) в этом представлении есть полином от самосопряжённого оператора, аппроксимирующий его унитарную экспоненту с ошибкой (7). Матрица $H$ получается из $\sum z_dL_d$ умножением на перестановку $e_c\mapsto e_{Sc}$ и ортогональным сжатием. Поэтому $\|H\|\le1+\eps$.

При каждом перемножении двоичных степеней скалярное произведение вычисляется точно в масштабе $Q^2$. Каждая компонента округляется к ближайшему целому в масштабе $Q$, затем элементы с обеими компонентами по модулю $\le6$ зануляются. Ошибка элемента не превосходит $6.5\sqrt2/Q<10/Q$. На блоке размера $d$ ошибка в операторной норме не больше $\eta_d=10d/Q$. Для дерева произведения с $j$ листьями вычисленная норма не выше $B^{2j-1}$, где $B=1+\eps+\eta_d$. Действительно, при слиянии поддеревьев верхняя граница до округления есть $B^{2j-2}$ и $B^{2j-2}(B-1)\ge\eta_d$. Индукция по тому же дереву даёт ошибку относительно $H^j$ не выше $(j-1)\eta_d B^{2j-2}$: в разности двух произведений заменяем сначала левый, затем правый множитель и прибавляем ошибку последнего округления. Повторное использование одной степени допускает разворачивание в дерево. При $j\le209$ нормы меньше $1.000001$. Неравенство Коши для строки и столбца ограничивает сумму модулей точных произведений числом $1.000001^2Q^2$, что помещается в знаковые 128 бит, в том числе для частичных сумм действительной и мнимой частей.

Обозначим $z=\Tr U_N^{209}/(N\mu)$. Полученный центр округлённо равен $1.82819\cdot10^{-5}+1.02319\cdot10^{-6}i$. Для полного рационального центра $z_0$ из архива
\[
 |z-z_0|\le1.000001\left(209\eps+209\sum_{d}\frac{10d^2}{Q}+10^{-60}\right)
 <2.1\cdot10^{-7}.                                                        \tag{8}
\]
Первое слагаемое взято до переноса по (4), поэтому множитель $d$ в нём отсутствует. Второе учитывает $|\Tr E|\le d\|E\|$ на четырёх блоках, третье покрывает замену остаточной непрерывной фазы. Здесь точный центр $z_0=(42155281424489+2359315878140i)/2^{61}$. Получается $\Re z>1.8\cdot10^{-5}$, что доказывает (1).
\newpage
\sect{5. Что обнаруживает один большой след}
Пишем $p_{m,N}=\Tr U_N^m$, $X_L=\Xi_N([0,L])$. Для $0<L<N$ и любого целого $m\ne0$ интегрирование по сдвигу даёт
\[
 \left|\EE_\alpha e^{2\pi i m\alpha}(X_L/L-1)\right|
 =\frac{|p_{m,N}|}{N}\left|\frac{\sin(\pi m L/N)}{\pi m L/N}\right|.
\]
Так как $\EE(X_L/L-1)=0$, получаем для $G_L=(1-X_L/L)_+\in[0,1]$
\[
 \EE G_L=\tfrac12\EE|X_L/L-1|\ge\tfrac{|p_{m,N}|}{2N}
 |\sinc(\pi m L/N)|,\qquad \sinc u=\sin(u)/u .                          \tag{9}
\]
У $\si$ из $\rho_2=1-|K|^2$ следует $\operatorname{Var}X_L
=L-\int_{[0,L]^2}|K(x,y)|^2dx\,dy\le L$, откуда $\EE_{\si}G_L\le1/(2\sqrt L)$. При $L=10^{12}$, $m=209$, $N=h_{209}$ формулы (1), (9) и $\sinc y\ge1-y^2/6$ в данном диапазоне дают
\[
 \EE_\alpha G_L-\EE_{\si}G_L>8.4\cdot10^{-6}.                          \tag{10}
\]
Границы интервала почти наверное не содержат точек $\si$, поэтому $G_L$ годится для проверки слабого предела. Для абелева теста $L_N(s)=N^{-2}\sum_{m\ge1}e^{-2sm/N}|p_{m,N}|^2$ достаточно сохранить один член. При $s=10^6$ из $e^{-y}\ge1-y$ следует
\[
 L_{h_{209}}(s)-\frac{1-e^{-2s}}{4s^2}
 \ge (1-418s/N)(1.8\cdot10^{-5})^2-\frac1{4s^2}>3.2\cdot10^{-10}.
\]

\subt{Непрерывный тест меньшей поддержки} Для $T=10^6$ возьмём $f_T(y)=(1-|y|/T)_+$ и $G=(1-\Xi_N(f_T)/T)_+$. Вместо множителя в (9) возникает $\sinc^2(\pi mT/N)$, поскольку $\widehat f_T(t)=T\sinc^2(\pi Tt)$. У $\si$ дисперсия этой линейной статистики до деления на $T$ равна $\int |\widehat f_T(t)|^2\min(|t|,1)dt$. Это следует из формулы для $\rho_2$ и $\widehat{|K(\cdot,0)|^2}=(1-|t|)_+$. Дисперсия ограничена интегралом $\int |u|\sinc^4(\pi u)du\le 2/\pi^2<1/4$; достаточно разделить интегрирование при $|u|=1/\pi$ и использовать $|\sin y|\le\min(|y|,1)$. Поэтому $\EE_{\si}G<1/(4T)$, а
\[
 \EE_\alpha G-\EE_{\si}G>8.74\cdot10^{-6}\qquad(N=h_{209}).
\]

\subt{Необходимое условие непосредственно для слабого закона} Если $\Xi_N\Rightarrow\si$, то для каждого $C<\infty$ с $C>0$
\[
 \max_{1\le m\le CN}|p_{m,N}|/N\longrightarrow0 .                       \tag{11}
\]
Допустим противное, выберем подпоследовательность с $m/N\to t\in[0,C]$ и следом, отделённым от нуля. Для фиксированного $L$ рассмотрим
$Z_N=L^{-1}\int_0^L e^{2\pi i(m/N)y}(d\Xi_N(y)-dy)$. Тогда $|p_{m,N}|/N\le\EE|Z_N|$ взятием коэффициента Фурье по $\alpha$ как в (9). Модуляции равномерно сходятся на окне к $e^{2\pi i t y}$; ошибка в $L^1$ ограничивается удвоенной равномерной ошибкой модуляции. Счётчик $X_L$ сходится по распределению и имеет ту же среднюю $L$, что и предел. Поэтому семейство $X_L$ равномерно интегрируемо: применим сходимость к $X_L\wedge B$ и затем увеличим $B$. Значит, $\EE|Z_N|$ сходится к своему аналогу у $\si$. Для $f=\mathbf1_{[0,L]}e^{2\pi i ty}$ дисперсия $\int f\,d\Xi$ не превосходит $\int |f|^2=L$, поскольку вычитаемая квадратичная форма с ядром $|K(x,y)|^2$ положительна (её фурье-множитель $(1-|u|)_+$). Итак, верхний предел нормированного следа не превосходит $1/\sqrt L$. Устремление $L$ к бесконечности доказывает (11).
\newpage
\sect{6. Слабый закон и парные тесты}
Введём положительную меру разностей фаз, включая совпадения,
\[
 \gamma_N=N^{-1}\sum_{j,l=1}^N\sum_{n\in\Z}\delta_{N(\theta_j-\theta_l+n)},
 \quad
 \gamma_{\rm s}=\delta_0+(1-|K(y,0)|^2)\,dy.
\]
Для $b\ge0$ из исходного класса положим $g_b(y)=\int b(x)b(x+y)dx$ и $h_s(y)=s/(s^2+\pi^2y^2)$. При обозначениях исходной задачи
\[
 2\Delta_{2,N}(b)=(\gamma_N-\EE_{\rm CUE}\gamma_N)(g_b),
 \qquad 1+2 L_N(s)=\gamma_N(h_s).                                        \tag{12}
\]
В первой формуле разложение произведения выделяет одинаковые индексы с суммарным вкладом $\sum_n g_b(Nn)$, общим для обоих спектров. Во второй формуле суммирование Пуассона даёт коэффициенты Фурье $e^{-2s|m|/N}/N$ у периодизации $h_s$ на окружности длины $N$. Тем самым нулевой член равен $1$. У эталона $\EE_{\rm CUE}\gamma_N(g_b)\to\gamma_{\rm s}(g_b)$: для $m\ge1$ коэффициент квадрата следа равен $\min(m,N)$, и возникает риманова сумма для $|\widehat b(t)|^2\min(|t|,1)$. Эквивалентно, $\widehat\gamma_{\rm s}=\delta_0+\min(|t|,1)dt$ с конвенцией $e^{-2\pi i ty}$. Также $\gamma_{\rm s}(h_s)=1+(1-e^{-2s})/(2s^2)$.

\subt{Односторонний тест} Из слабой сходимости $\Xi_N\Rightarrow\si$ следует
\[
 \liminf_N\gamma_N(g)\ge\gamma_{\rm s}(g)
 \qquad(g\ge0\text{ непрерывна}).                                        \tag{13}
\]
Возьмём $\rho\ge0$ непрерывную компактной поддержки с интегралом $1$. Для $g\in C_c$ левая статистика есть ожидание
$\iint\rho(x)g(y-x)d\Xi_N(x)d\Xi_N(y)$ по стационарности и явной периодизации. Функционал непрерывен в vague-топологии: на компактном прямоугольнике ядро равномерно приближается суммами произведений непрерывных тестов, а массы компактных множеств ограничены на сходящейся последовательности мер. Усечение функционала на уровне $B$, затем $B\uparrow\infty$, даёт (13). Монотонное исчерпание расширяет вывод на некомпактные $g$. Следовательно, отрицательный предельный дефицит парного теста из (12) опроверг бы BGS-F даже без предположений о вторых моментах.

\subt{Один парный тест при известном слабом законе} Если одновременно имеют место BGS-F и (P), то $\gamma_N\to\gamma_{\rm s}$ на всех непрерывных тестах $f$ с $\sum_{n\in\Z}\sup_{[n,n+1]}|f|<\infty$. В частности, следуют (A), (R) и вся вторая строка иерархии (H), то есть каждый допустимый $b$ при $k=2$. Доказательство завершается на следующей странице.
\newpage
\sect{7. Доказательство парного утверждения и граница результата}
В (P) берётся $b_*(x)=|\mathcal F^{-1}\phi(x)|^2/(2I^2)$, где $I=\int\phi$ и $\phi(t)=e^{-1/(1-t^2)}$ при $|t|<1$, иначе $\phi=0$. Функция $g_*=g_{b_*}$ строго положительна на всей оси. Действительно, $\mathcal F^{-1}\phi$ продолжается в целую ненулевую функцию, поэтому вещественные нули $b_*$ изолированы и произведение под интегралом для $g_*$ положительно почти всюду. По (12) условие (P) даёт $\gamma_N(g_*)\to\gamma_{\rm s}(g_*)$. На каждом компакте минимум $g_*$ положителен, откуда следует относительная vague-компактность $\gamma_N$.

Пусть $\nu$ --- частичный предел. Из (13) следует $\nu\ge\gamma_{\rm s}$. Усечение $g_*$ компактными срезками даёт
$\nu(g_*)\le\liminf\gamma_N(g_*)=\gamma_{\rm s}(g_*)$. Положительная мера $\nu-\gamma_{\rm s}$ интегрирует всюду положительную функцию в ноль и потому равна нулю. Единственность частичного предела даёт vague-сходимость.

Для получения сходимости на заявленном классе остаются хвосты. Положим $h(y)=(1-|y|)_+$, $Y_t=\Xi_N([t,t+1])$. Интегрирование по общему сдвигу даёт $(h*\gamma_N)(t)=\EE Y_0Y_t\le\EE Y_0^2=(h*\gamma_N)(0)$ по Коши и стационарности. Правая часть равномерно ограничена контролем $\gamma_N([-1,1])$. Так как $h\ge1/2$ на $[-1/2,1/2]$, массы всех единичных интервалов равномерно ограничены по месту и по $N$. Теперь ряд супремумов мажорирует хвосты, что завершает доказательство. Для (R) используется точная связь (8.4)--(8.5) исходной постановки.

\subt{Независимая проверка переноса} Для $k=3,5,7$ ($N=7,41,239$) прямое возведение матрицы в степень сравнено с двумя вычислениями: перебором упорядоченных слов произвольного пятичленного полинома и формулой (4) для $V=I+(3i/8)A$. Во втором случае $A$ независимо собрано как трёхдиагональная циклическая матрица: правый внедиагональный элемент равен $e^{-i\alpha_0}r((x+\bar2)/N)/200$, левый задаётся сопряжением. Максимальное расхождение двух значений нормированного следа во втором тесте меньше $7\cdot10^{-16}$. Эти сравнения проверяют фазовую конвенцию; в (8) используются целочисленные операции.

\subt{Что остаётся для решения} На пеллевых размерностях $k\sim\log(2N)/\log(1+\sqrt2)$. Равномерное продолжение (1) с положительной константой на бесконечно много нечётных $k$ противоречило бы (11) и опровергло бы исходную слабую сходимость. При фиксированной $\tau$ параметр $\xi=\tau h_k$ неограниченно растёт. Константы обрезки из раздела 3 получены только для $k=209$. Тождество (4) сохраняет силу, но оценки его следа снизу в таком режиме пока нет. И само исчезновение этого пика было бы лишь одним необходимым тестом. Для положительного решения по-прежнему требуется контроль локальных функционалов или полной знаковой иерархии на гейзенберговских временах.

\subt{Источники и воспроизведение} Определения и обозначения (P), (A), (R), (H) взяты из приложенной постановки BGS-F от 19.11.2026, ревизия \texttt{e2757e4}. О коротких периодах линейных кошачьих карт см. \href{https://doi.org/10.1007/s00220-003-0888-3}{F. Faure, S. Nonnenmacher, S. De Bièvre (2003)}. О детерминантных процессах см. \href{https://arxiv.org/abs/math/0002099}{A. Soshnikov (2000)}. Архив содержит исходную постановку, текст отчёта, рациональные коэффициенты в форме алгоритма, матрицы, результаты целочисленного расчёта, независимую маломерную проверку и команды для воспроизведения.
\end{document}
